
isolée sont des invariants géométriques tout à fait non triviaux, et qui auront sans doute leur rôle à jouer dans l'étude ``arithmétique'' de ces anneaux locaux.

Bien entendu, une fois en possession d'une notion raisonnable d'anneau de Chow pour un schéma noethérien régulier admettant un Module inversible ample, il y a lieu de généraliser à ce cas les relations établies resp. conjecturées au n\textsuperscript{o} 4 entre cet anneau et l'anneau @@M0@@, et à formuler en termes de cet anneau la question de la ``formule de Riemann--Roch sans dénominateurs'' énoncée au n\textsuperscript{o} 3.

\subsection*{BIBLIOGRAPHIE}
\begin{enumerate}[label={[\arabic*]}]
\item M. F. Atiyah et F. Hirzebruch, The Riemann--Roch theorem for analytic embeddings, Topology, vol. 1, p. 151-166 (1962).
\item A. Borel et J.P. Serre, Le théorème de Riemann--Roch, Bull. Soc. Math. France, vol. 86, p. 97-136 (1958).
\item R. Bott et H. Samelson, Amer. Journ. of Math. vol 80, p. 1004 (1958).
\item C. Chevalley, Les classes d'équivalence rationnelle, I et II, Séminaire à l'ENS 1958, exposés 2 et 3.
\item C. Chevalley, Classification des groupes de Lie algébriques, Séminaire à l'ENS 1956/58.
\item A. Dold et D. Puppe, Homologie nicht-additiver Funktoren. Anwendungen, Annales de l'Inst. Fourier t. XI, p. 201-312 (1961).
\item H. Grauert, Ein Theorem der analytischen Garbentheorie, Pub. Math. n\textsuperscript{o} 5, p. 5-64 (1960).
\item A. Grothendieck, Classes de faisceaux et théorème de Riemann--Roch, notes multigraphiées par les soins de l'Institute for Advanced Study, 1957 (travail cité [RRR] dans le présent exposé, et qui est reproduit dans notre séminaire sous forme d'Appendice au Chap. O).
\item A. Grothendieck, La théorie des classes de Chern, Bull. Soc. Math. France vol. 86, p. 137-154 (1958).
